\section{Solarity and Impassibility}

\begin{quotex}
Le cœur a ses raisons, que la raison ne connait pas. \flright{\textsc{Blaise Pascal}}

And you shall know the truth, and the truth shall make you free. \flright{\textsc{John 8:32}}

\end{quotex}
Many people expect to experience physical effects and/or emotional consolations from the meditation and prayer practice. I'm not denying that effects like visions, auditory phenomena, and condolences are possible. Rather they are not the goal. A fortiori, they usually diminish as you advance and it would be a mistake to expect them all the time.

At a recent Sunday sermon, the priest brought up the same point. He said that because God is spirit, He has no emotions, at least not as we understand them. Hence, we should not seek such emotional consolations. Of course, he was referring to God's impassibility but probably didn't want to explain that teaching to random lay people.

\paragraph{Impassibility}
Impassibility was one of the preternatural gifts before the Fall. (The others are integrity, immortality, and infused knowledge or gnosis.) If the aim of the spiritual life is theosis, to become more like God, then we need to recover the quality of impassibility that was lost. As this applies to God, the philosopher Edward Feser explains\footnote{\url{https://edwardfeser.blogspot.com/2014/09/olson-contra-classical-theism.html}}: 

\begin{quotex}
the claim that classical theism makes God out to be ``unemotional" is ambiguous. If by an ``emotion" we mean a state that comes upon us episodically, that varies in its intensity, that has physiological aspects like increased heart rate and bodily sensations, etc., then it is certainly true that the classical theist maintains that God cannot possibly have such states. However, if the insinuation is that classical theism makes God out to be ``unemotional" in a way that entails that he cannot be said to love us, to be angry at sin, etc., then that is certainly false. To love is to will the good of another, and for the classical theist God certainly wills our good, acts providentially so that we attain what is good for us, etc. Hence he loves us. The classical theist also holds that God wills that sin be punished, and acts so that those who are unrepentant are in fact punished. Hence he is in that sense wrathful at sin. And so forth. Hardly ``apathetic." 

\end{quotex}
Our impassibility is analogous to God's, but not identical. Obviously, as embodied beings, we do indeed experience emotions that spontaneously arise in us and with intense bodily sensations. Nevertheless, we can achieve some freedom from negative emotions like worry, doubt, fear, hate, apathy and despair, all of which imprison us to some degree. Of course, unlike in Eden, we are subject still to irreprehensible qualities such as hunger, thirst, pain, weariness, death, disease.

\paragraph{Practical Esoterism}
\textbf{Valentin Tomberg} has provided us a different standard by which to measure spiritual progress: viz., the sequence mysticism, gnosis, magic, philosophy. These can be briefly summarized.

\begin{itemize}
\item \textbf{Mysticism}: The first step is inner contact with the divine, beyond created being. This requires the body to be relaxed and the soul to be serene. This means that inner serenity is the starting point rather than the result. 
\item \textbf{Gnosis}: Gnosis is mysticism that has become conscious of itself, i.e., it is mystical experience that has been transformed into higher knowledge. This is direct, unmediated intuition, not a theory. 
\item \textbf{Magic}: Magic is the power of the invisible and spiritual over the visible and material. The aim of magic is ``liberating action, i.e. the restoration of freedom to beings who have partially or totally lost it", as Tomberg phrases it. Hence, gnosis is knowing the truth, which then sets us free. 
\item \textbf{Philosophy}: Esoteric philosophy is the culmination of the sequence. Its fundamental requirement is not to substitute one intellectual conception of the world with another, but rather to create in the individual a new ``dimension" and a new depth of life. 
\end{itemize}
Note that this does not include emotional consolation nor any physical sensations as necessary aspects of the mystical experience. To the contrary, emotions need to be silenced. Tomberg warns us of the danger of not following through the whole sequence.

\begin{quotex}
Mysticism which has not given birth to gnosis, magic and Hermetic philosophy—such a mysticism must, sooner or later, necessarily degenerate into spiritual enjoyment" or ``intoxication". The mystic who wants only the experience of mystical states without understanding them, without drawing practical conclusions from them for life, and without wanting to be useful to others, who forgets everyone and everything in order to enjoy the mystical experience, can be compared to a spiritual drunkard. 

\end{quotex}
\paragraph{Three Vows}
To prepare us, Tomberg provides us with the esoteric understanding of the three traditional vows. These are suitable not just to the religious, but especially to those living in the world.

\begin{itemize}
\item \textbf{Obedience}: The vow of obedience is the practice of silencing personal desires, emotions and imagination in the face of reason and conscience. 
\item \textbf{Poverty}: The vow of poverty is the practice of inner emptiness, which is established as a consequence of the silence of personal desires, emotions and imagination so that the soul is capable of receiving from above the revelation of the word, the life and the light. 
\item \textbf{Chastity}: The vow of chastity means to say the putting into practice of the resolution to live according to solar law, without covetousness and without indifference. 
\end{itemize}
In other words, obedience means that personal desires, emotions, and fantasies must be made subservient to the intellect and the moral conscience. Spiritual poverty means the silencing of such desires, emotions, and fantasies.

\paragraph{Solarity}
Tomberg then explains the meaning of the solar law, or solarity.

\begin{quotex}
For virtue is boring and vice is disgusting. But that which comes from the depths of the heart is neither boring nor disgusting. The foundation of the heart is love. The heart lives only when it loves. It is then like the sun. And chastity is the state of the human being in which the heart, having become solar, is the centre of gravity. 

\end{quotex}
This describes the awakening of the heart chakra, as he explains:

\begin{quotex}
Chastity is the state of the human being where the twelve-petalled lotus (anahata in Indian esoterism) is awakened and becomes the sun of the microcosmic ``planetary system". The three lotus-centres situated below it (the ten-petalled, the six-petalled, and the four-petalled) begin then to function in conformity with the life of the heart (the twelve-petalled lotus), i.e. ``according to solar law". When they do this, the person is chaste, no matter whether he or she is celibate or married. 

\end{quotex}
The three lowest chakras are the seats of fear, sex, and hunger. When these chakras are subservient to the solar law, they are balanced and beneficial. Otherwise, they are disordered, abnormal, and come to dominate consciousness when they should be in line with intelligence and morality.

For more on the chakras as described by Johann Gichtel and Oscar Hinze, see \textit{The Degrees of Existence}\footnote{\url{http://www.gornahoor.net/?p=8618}}.

For the lower chakras, see \textit{Sex and Tropism}\footnote{\url{https://www.gornahoor.net/?p=8349}}.



\flrightit{Posted on 2018-12-21 by Cologero }
